\documentclass[10pt,a4paper]{amsart}
\usepackage{amsmath,amssymb,amsthm,graphicx}
\usepackage{verbatim}
\newenvironment{reference}{\comment}{\endcomment}
%\newenvironment{reference}{}{}
\newenvironment{slogan}{\comment}{\endcomment}
\newenvironment{history}{\comment}{\endcomment}
% For commutative diagrams we use Xy-pic
\usepackage[all]{xy}

% We use 2cell for 2-commutative diagrams.
\xyoption{2cell}
% Theorem environments.
%
\theoremstyle{plain}
\newtheorem{theorem}[subsection]{Theorem}
\newtheorem{proposition}[subsection]{Proposition}
\newtheorem{lemma}[subsection]{Lemma}

\theoremstyle{definition}
\newtheorem{definition}[subsection]{Definition}
\newtheorem{example}[subsection]{Example}
\newtheorem{exercise}[subsection]{Exercise}
\newtheorem{situation}[subsection]{Situation}

\theoremstyle{remark}
\newtheorem{remark}[subsection]{Remark}
\newtheorem{remarks}[subsection]{Remarks}

\numberwithin{equation}{subsection}

% Macros
%
\def\lim{\mathop{\mathrm{lim}}\nolimits}
\def\colim{\mathop{\mathrm{colim}}\nolimits}
\def\Spec{\mathop{\mathrm{Spec}}}
\def\Hom{\mathop{\mathrm{Hom}}\nolimits}
\def\Ext{\mathop{\mathrm{Ext}}\nolimits}
\def\SheafHom{\mathop{\mathcal{H}\!\mathit{om}}\nolimits}
\def\SheafExt{\mathop{\mathcal{E}\!\mathit{xt}}\nolimits}
\def\Sch{\mathit{Sch}}
\def\Mor{\mathop{\mathrm{Mor}}\nolimits}
\def\Ob{\mathop{\mathrm{Ob}}\nolimits}
\def\Sh{\mathop{\mathit{Sh}}\nolimits}
\def\NL{\mathop{N\!L}\nolimits}
\def\CH{\mathop{\mathrm{CH}}\nolimits}
\def\proetale{{pro\text{-}\acute{e}tale}}
\def\etale{{\acute{e}tale}}
\def\QCoh{\mathit{QCoh}}
\def\Ker{\mathop{\mathrm{Ker}}}
\def\Im{\mathop{\mathrm{Im}}}
\def\Coker{\mathop{\mathrm{Coker}}}
\def\Coim{\mathop{\mathrm{Coim}}}

% Boxtimes
%
\DeclareMathSymbol{\boxtimes}{\mathbin}{AMSa}{"02}

%
% Macros for moduli stacks/spaces
%
\def\QCohstack{\mathcal{QC}\!\mathit{oh}}
\def\Cohstack{\mathcal{C}\!\mathit{oh}}
\def\Spacesstack{\mathcal{S}\!\mathit{paces}}
\def\Quotfunctor{\mathrm{Quot}}
\def\Hilbfunctor{\mathrm{Hilb}}
\def\Curvesstack{\mathcal{C}\!\mathit{urves}}
\def\Polarizedstack{\mathcal{P}\!\mathit{olarized}}
\def\Complexesstack{\mathcal{C}\!\mathit{omplexes}}
% \Pic is the operator that assigns to X its picard group, usage \Pic(X)
% \Picardstack_{X/B} denotes the Picard stack of X over B
% \Picardfunctor_{X/B} denotes the Picard functor of X over B
\def\Pic{\mathop{\mathrm{Pic}}\nolimits}
\def\Picardstack{\mathcal{P}\!\mathit{ic}}
\def\Picardfunctor{\mathrm{Pic}}
\def\Deformationcategory{\mathcal{D}\!\mathit{ef}}
\usepackage{cleveref}
\crefname{originallemma}{original lemma}{original lemmas}
\Crefname{originallemma}{Original lemma}{Original lemmas}
\makeatletter
\newlabel{lemma-base-change-finite-type}{{lemma-base-change-finite-type}{source}{}{sourcepointer-1}{}}
\newlabel{lemma-base-change-finite-type@cref}{{[originallemma][0][]lemma-base-change-finite-type}{[1][0][]source}{}{}{}}
\makeatother
\crefname{originallemma}{original lemma}{original lemmas}
\Crefname{originallemma}{Original lemma}{Original lemmas}
\makeatletter
\newlabel{lemma-composition-finite-presentation}{{lemma-composition-finite-presentation}{source}{}{sourcepointer-2}{}}
\newlabel{lemma-composition-finite-presentation@cref}{{[originallemma][0][]lemma-composition-finite-presentation}{[1][0][]source}{}{}{}}
\makeatother
\crefname{originallemma}{original lemma}{original lemmas}
\Crefname{originallemma}{Original lemma}{Original lemmas}
\makeatletter
\newlabel{lemma-composition-flat}{{lemma-composition-flat}{source}{}{sourcepointer-3}{}}
\newlabel{lemma-composition-flat@cref}{{[originallemma][0][]lemma-composition-flat}{[1][0][]source}{}{}{}}
\makeatother
\crefname{originallemma}{original lemma}{original lemmas}
\Crefname{originallemma}{Original lemma}{Original lemmas}
\makeatletter
\newlabel{lemma-descent-finite-presentation}{{lemma-descent-finite-presentation}{source}{}{sourcepointer-4}{}}
\newlabel{lemma-descent-finite-presentation@cref}{{[originallemma][0][]lemma-descent-finite-presentation}{[1][0][]source}{}{}{}}
\makeatother
\crefname{originallemma}{original lemma}{original lemmas}
\Crefname{originallemma}{Original lemma}{Original lemmas}
\makeatletter
\newlabel{lemma-descent-flat}{{lemma-descent-flat}{source}{}{sourcepointer-5}{}}
\newlabel{lemma-descent-flat@cref}{{[originallemma][0][]lemma-descent-flat}{[1][0][]source}{}{}{}}
\makeatother
\crefname{originallemma}{original lemma}{original lemmas}
\Crefname{originallemma}{Original lemma}{Original lemmas}
\makeatletter
\newlabel{lemma-gerbe-fppf}{{lemma-gerbe-fppf}{source}{}{sourcepointer-6}{}}
\newlabel{lemma-gerbe-fppf@cref}{{[originallemma][0][]lemma-gerbe-fppf}{[1][0][]source}{}{}{}}
\makeatother
\crefname{originallemma}{original lemma}{original lemmas}
\Crefname{originallemma}{Original lemma}{Original lemmas}
\makeatletter
\newlabel{lemma-inertia}{{lemma-inertia}{source}{}{sourcepointer-7}{}}
\newlabel{lemma-inertia@cref}{{[originallemma][0][]lemma-inertia}{[1][0][]source}{}{}{}}
\makeatother
\crefname{originalproposition}{original proposition}{original propositions}
\Crefname{originalproposition}{Original proposition}{Original propositions}
\makeatletter
\newlabel{proposition-when-gerbe}{{proposition-when-gerbe}{source}{}{sourcepointer-8}{}}
\newlabel{proposition-when-gerbe@cref}{{[originalproposition][0][]proposition-when-gerbe}{[1][0][]source}{}{}{}}
\makeatother
\crefname{originalreference}{original reference}{original references}
\Crefname{originalreference}{Original reference}{Original references}
\makeatletter
\newlabel{spaces-morphisms-lemma-noetherian-finite-type-finite-presentation}{{spaces-morphisms-lemma-noetherian-finite-type-finite-presentation}{source}{}{sourcepointer-9}{}}
\newlabel{spaces-morphisms-lemma-noetherian-finite-type-finite-presentation@cref}{{[originalreference][0][]spaces-morphisms-lemma-noetherian-finite-type-finite-presentation}{[1][0][]source}{}{}{}}
\makeatother
\crefname{originalreference}{original reference}{original references}
\Crefname{originalreference}{Original reference}{Original references}
\makeatletter
\newlabel{stacks-properties-lemma-composition-surjective}{{stacks-properties-lemma-composition-surjective}{source}{}{sourcepointer-10}{}}
\newlabel{stacks-properties-lemma-composition-surjective@cref}{{[originalreference][0][]stacks-properties-lemma-composition-surjective}{[1][0][]source}{}{}{}}
\makeatother
\crefname{originalreference}{original reference}{original references}
\Crefname{originalreference}{Original reference}{Original references}
\makeatletter
\newlabel{stacks-properties-lemma-unique-point-better}{{stacks-properties-lemma-unique-point-better}{source}{}{sourcepointer-11}{}}
\newlabel{stacks-properties-lemma-unique-point-better@cref}{{[originalreference][0][]stacks-properties-lemma-unique-point-better}{[1][0][]source}{}{}{}}
\makeatother
\title{A locally Noetherian reduced stack with a single point is a gerbe over a field}
\author{The Stacks Project Authors}
\date{stacks-project:stacks-morphisms.tex:local-snapshot-20261012; GNU FDL 1.2 or later}
\begin{document}
\maketitle
\noindent\textit{Editorial note.} All delivered mathematical text is the author's own, copied from the pinned Stacks Project chapter source sources/stacks-algebraic-stacks-morphisms/pinned-original-source.tex; no statement or proof step was rewritten, reordered or removed.\par\medskip
\noindent\textit{Editorial note.} Preamble provenance: the Stacks Project keeps one single shared preamble for every chapter; this card takes it directly from the staged sources/stacks-algebraic-stacks-morphisms/preamble.tex as three verbatim spans (lines 8-12, 14-18 and 154-217), so no derived file was written and nothing was added or redefined. The wrapper class and the base packages are supplied by the assembler.\par\medskip
\noindent\textit{Editorial note.} Bibliography provenance: neither the counted statement nor its proof contains a citation (this lane rescans the whole span before writing the card), so the card carries no bibliography and no entry was reconstructed.\par\medskip
\noindent\textit{Editorial note.} Selection provenance: this card was opened under the file-based control contract (CONTROL.md), with the mechanical pre-card gate run on its statement and proof interval before the card was written.\par\medskip
\noindent\textit{Allowed standard prerequisites (editorial).}\begin{itemize}
\item The existence of a surjective flat locally finitely presented morphism from the spectrum of a field onto a stack whose space is a single point. Invoked at the start of the counted proof (source lines 6580-6583) as a cross-chapter pointer. Produces the field-valued cover used to test the inertia.
\item The inertia and its base change properties. Invoked inside the counted proof (source lines 6584-6592). Shows that the inertia over the field is locally of finite type, hence of finite presentation.
\item The descent statements for flatness and local finite presentation. Invoked inside the counted proof (source lines 6593-6596). Descends the properties of the inertia to the stack.
\item The characterization of gerbes by flatness and local finite presentation of the inertia, and the surjectivity flatness and finite presentation of a gerbe morphism. Invoked inside the counted proof (source lines 6597-6608). Concludes that the stack is a gerbe and that its base has a single point and is locally Noetherian and reduced.
\end{itemize}\medskip
\ifdefined\hypertarget\hypertarget{sourcepointer-1}{}\fi
\noindent\textit{Source reference.} Original lemma lemma-base-change-finite-type: sources/stacks-algebraic-stacks-morphisms/pinned-original-source.tex, source lines 3501--3501. Navigation pointer only: the counted text refers to the Stacks item lemma-base-change-finite-type, which is an existing labeled result of this same chapter and is not part of this card. The pointer carries the target's own label as its number and its definition line as its span; its content is not reproduced here.\par\medskip
\ifdefined\hypertarget\hypertarget{sourcepointer-2}{}\fi
\noindent\textit{Source reference.} Original lemma lemma-composition-finite-presentation: sources/stacks-algebraic-stacks-morphisms/pinned-original-source.tex, source lines 5614--5614. Navigation pointer only: the counted text refers to the Stacks item lemma-composition-finite-presentation, which is an existing labeled result of this same chapter and is not part of this card. The pointer carries the target's own label as its number and its definition line as its span; its content is not reproduced here.\par\medskip
\ifdefined\hypertarget\hypertarget{sourcepointer-3}{}\fi
\noindent\textit{Source reference.} Original lemma lemma-composition-flat: sources/stacks-algebraic-stacks-morphisms/pinned-original-source.tex, source lines 5363--5363. Navigation pointer only: the counted text refers to the Stacks item lemma-composition-flat, which is an existing labeled result of this same chapter and is not part of this card. The pointer carries the target's own label as its number and its definition line as its span; its content is not reproduced here.\par\medskip
\ifdefined\hypertarget\hypertarget{sourcepointer-4}{}\fi
\noindent\textit{Source reference.} Original lemma lemma-descent-finite-presentation: sources/stacks-algebraic-stacks-morphisms/pinned-original-source.tex, source lines 5833--5833. Navigation pointer only: the counted text refers to the Stacks item lemma-descent-finite-presentation, which is an existing labeled result of this same chapter and is not part of this card. The pointer carries the target's own label as its number and its definition line as its span; its content is not reproduced here.\par\medskip
\ifdefined\hypertarget\hypertarget{sourcepointer-5}{}\fi
\noindent\textit{Source reference.} Original lemma lemma-descent-flat: sources/stacks-algebraic-stacks-morphisms/pinned-original-source.tex, source lines 5389--5389. Navigation pointer only: the counted text refers to the Stacks item lemma-descent-flat, which is an existing labeled result of this same chapter and is not part of this card. The pointer carries the target's own label as its number and its definition line as its span; its content is not reproduced here.\par\medskip
\ifdefined\hypertarget\hypertarget{sourcepointer-6}{}\fi
\noindent\textit{Source reference.} Original lemma lemma-gerbe-fppf: sources/stacks-algebraic-stacks-morphisms/pinned-original-source.tex, source lines 6287--6287. Navigation pointer only: the counted text refers to the Stacks item lemma-gerbe-fppf, which is an existing labeled result of this same chapter and is not part of this card. The pointer carries the target's own label as its number and its definition line as its span; its content is not reproduced here.\par\medskip
\ifdefined\hypertarget\hypertarget{sourcepointer-7}{}\fi
\noindent\textit{Source reference.} Original lemma lemma-inertia: sources/stacks-algebraic-stacks-morphisms/pinned-original-source.tex, source lines 1220--1220. Navigation pointer only: the counted text refers to the Stacks item lemma-inertia, which is an existing labeled result of this same chapter and is not part of this card. The pointer carries the target's own label as its number and its definition line as its span; its content is not reproduced here.\par\medskip
\ifdefined\hypertarget\hypertarget{sourcepointer-8}{}\fi
\noindent\textit{Source reference.} Original proposition proposition-when-gerbe: sources/stacks-algebraic-stacks-morphisms/pinned-original-source.tex, source lines 6319--6319. Navigation pointer only: the counted text refers to the Stacks item proposition-when-gerbe, which is an existing labeled result of this same chapter and is not part of this card. The pointer carries the target's own label as its number and its definition line as its span; its content is not reproduced here.\par\medskip
\ifdefined\hypertarget\hypertarget{sourcepointer-9}{}\fi
\noindent\textit{Source reference.} Original reference spaces-morphisms-lemma-noetherian-finite-type-finite-presentation: sources/stacks-algebraic-stacks-morphisms/pinned-original-source.tex, source lines 3847--3847. Navigation pointer only: the counted text refers to the Stacks item spaces-morphisms-lemma-noetherian-finite-type-finite-presentation, which belongs to another chapter of the project and is not part of this card. The pointer carries the target's own label as its number because this chapter's staged snapshot does not contain the other chapter's tag table; its content is not reproduced here.\par\medskip
\ifdefined\hypertarget\hypertarget{sourcepointer-10}{}\fi
\noindent\textit{Source reference.} Original reference stacks-properties-lemma-composition-surjective: sources/stacks-algebraic-stacks-morphisms/pinned-original-source.tex, source lines 5403--5403. Navigation pointer only: the counted text refers to the Stacks item stacks-properties-lemma-composition-surjective, which belongs to another chapter of the project and is not part of this card. The pointer carries the target's own label as its number because this chapter's staged snapshot does not contain the other chapter's tag table; its content is not reproduced here.\par\medskip
\ifdefined\hypertarget\hypertarget{sourcepointer-11}{}\fi
\noindent\textit{Source reference.} Original reference stacks-properties-lemma-unique-point-better: sources/stacks-algebraic-stacks-morphisms/pinned-original-source.tex, source lines 4414--4414. Navigation pointer only: the counted text refers to the Stacks item stacks-properties-lemma-unique-point-better, which belongs to another chapter of the project and is not part of this card. The pointer carries the target's own label as its number because this chapter's staged snapshot does not contain the other chapter's tag table; its content is not reproduced here.\par\medskip
\begin{lemma}
\label{lemma-noetherian-singleton-stack-gerbe}
Let $\mathcal{Z}$ be a reduced, locally Noetherian algebraic stack
such that $|\mathcal{Z}|$ is a singleton. Then $\mathcal{Z}$ is a gerbe
over a reduced, locally Noetherian algebraic space $Z$ with $|Z|$ a
singleton.
\end{lemma}

\begin{proof}
By
Properties of Stacks, Lemma \ref{stacks-properties-lemma-unique-point-better}
there exists a surjective, flat, locally finitely presented morphism
$\Spec(k) \to \mathcal{Z}$ where $k$ is a field.
Then $\mathcal{I}_Z \times_\mathcal{Z} \Spec(k) \to \Spec(k)$
is representable by algebraic spaces and
locally of finite type (as a base change of
$\mathcal{I}_\mathcal{Z} \to \mathcal{Z}$, see
Lemmas \ref{lemma-inertia} and \ref{lemma-base-change-finite-type}).
Therefore it is locally of finite presentation, see
Morphisms of Spaces, Lemma
\ref{spaces-morphisms-lemma-noetherian-finite-type-finite-presentation}.
Of course it is also flat as $k$ is a field. Hence we may apply
Lemmas \ref{lemma-descent-flat} and
\ref{lemma-descent-finite-presentation}
to see that $\mathcal{I}_\mathcal{Z} \to \mathcal{Z}$ is flat and
locally of finite presentation. We conclude that $\mathcal{Z}$
is a gerbe by
Proposition \ref{proposition-when-gerbe}.
Let $\pi : \mathcal{Z} \to Z$ be a morphism to an algebraic space such
that $\mathcal{Z}$ is a gerbe over $Z$. Then $\pi$ is surjective, flat, and
locally of finite presentation by
Lemma \ref{lemma-gerbe-fppf}.
Hence $\Spec(k) \to Z$ is surjective, flat, and locally of finite
presentation as a composition, see
Properties of Stacks,
Lemma \ref{stacks-properties-lemma-composition-surjective}
and
Lemmas \ref{lemma-composition-flat} and
\ref{lemma-composition-finite-presentation}.
Hence by
Properties of Stacks, Lemma \ref{stacks-properties-lemma-unique-point-better}
we see that $|Z|$ is a singleton and that $Z$ is locally Noetherian
and reduced.
\end{proof}

\end{document}
